\documentclass[10pt,a4paper]{amsart}
\usepackage[margin=19mm]{geometry}
\usepackage{amsmath,amssymb,mathtools}
\usepackage[scr=rsfs,cal=euler]{mathalfa}
\usepackage{xfrac}
\usepackage[unicode,hidelinks,bookmarks=false]{hyperref}
\newcommand{\ie}{i.e.\ }
\newcommand{\cf}{cf.\ }
\newcommand{\resp}{respectively\ }
\newcommand{\Subsection}{\subsection}
\providecommand{\nobreakdash}{\nobreak\hbox{-}}
\setlength{\emergencystretch}{2em}
\multlinegap0pt
\usepackage{bm}
\usepackage{bbm}
\usepackage{dsfont}
\usepackage{xspace}
\usepackage{cancel}
\usepackage{mathtools}
\usepackage{amsthm}
\makeatletter
\@ifundefined{theorem}{\newtheorem{theorem}{Theorem}}{}
\@ifundefined{proposition}{\newtheorem{proposition}[theorem]{\bf Proposition}}{}
\makeatother
\makeatletter
\def\FF{\mathbb{F}}
\def\field{K}
\expandafter\ifx\csname proofof\endcsname\relax
\newenvironment{proofof}[1]{\noindent{\it Proof of
#1.}}{\hfill$\square$\\\mbox{}}
\fi
\makeatother
\title[Separating polynomial invariants over non-closed fields of f]{Separating polynomial invariants over non-closed fields of finite abelian groups}
\author{M\'aty\'as Domokos}
\date{Source: arXiv:2506.22889v1 (CC BY 4.0)}
\begin{document}
\maketitle

\noindent\emph{Source and licence.} Separating polynomial invariants over non-closed fields of finite abelian groups. Version arXiv:2506.22889v1 (CC BY 4.0), distributed under the Creative Commons Attribution 4.0 International licence, as recorded in the arXiv licence metadata for this exact version and shown on the abstract page https://arxiv.org/abs/2506.22889v1. Full rights notice: licenses/arxiv-2506.22889-CC-BY-4.0.txt.\par\smallskip

\noindent\textit{Editorial scope.} Counted unit: the Proposition whose source label is \texttt{prop:irrep abelian} in the article (source0.tex lines 537--550, source label \texttt{prop:irrep abelian}), delivered together with the article's own complete original proof of it (source0.tex lines 552--609, one \texttt{proof} environment of 58 lines). No mathematics is written, repaired or re-derived: statement and proof are the article's own text line for line. Required context retained uncounted: none. Every definition, display and label of those lines is retained unchanged. Omitted from this package and not redistributed: 1--536, 551--551, 610--1136. Nothing inside a retained excerpt is omitted or altered: every retained body is delivered exactly as the article's lines are written, so no omission receipt of any kind accompanies this package. Citations: the retained excerpts carry no citation command at all, so no bibliography entry of the article is transcribed and no reference is redistributed. Reference adaptation: none was needed and none was made; every reference inside the delivered unit resolves inside it, so no unresolved reference or citation is produced. Printed slips kept literal: no printed word, symbol, hypothesis or number of the counted statement or of its proof was altered. Layout and class: the article's own class is \texttt{amsart} with packages including hyperref, geometry, comment, graphicx, amssymb; the wrapper is the lane's pinned \texttt{amsart} editorial wrapper at 10pt on A4, which restates the theorem-like environments the retained text needs and adds the article's own macro definitions that the retained text needs (\texttt{\textbackslash FF}, \texttt{\textbackslash field}), with extra wrapper packages (bm, bbm, dsfont, xspace, cancel, mathtools). The delivered numbering is the unit's own. Assets: no retained span contains a figure environment, a graphics-inclusion command, a PDF-page inclusion, a table or a TikZ picture; no image file is copied into this package, the package declares no asset dependency and no image file is redistributed with it. Licence: version arXiv:2506.22889v1 of this article is distributed under the Creative Commons Attribution 4.0 International licence, as recorded in the arXiv licence metadata for this exact version and shown on the abstract page https://arxiv.org/abs/2506.22889v1. A full rights notice is delivered at licenses/arxiv-2506.22889-CC-BY-4.0.txt. Characters: every retained excerpt is pure ASCII, so the delivered engine, XeLaTeX with its default Latin Modern text font, reports no missing character.
\begin{proposition} \label{prop:irrep abelian} 
\begin{itemize}
\item[(i)] The representation $(\field^{\ell_\chi},\rho_\chi)$ is isomorphic to 
$\bigoplus_{\psi\in \Gamma\cdot\chi} (\field,\psi)$, where $\Gamma\cdot \chi$ is the 
$\Gamma$-orbit of $\chi$. 
\item[(ii)] The representation $(\FF^{\ell_\chi},\rho_\chi)$ is irreducible for all $\chi\in \widehat G/\Gamma$. 
\item[(iii)] The regular representation $(\mathcal{F}(G,\FF),\rho_{\mathrm{reg}})$ of $G$ over $\FF$ 
is isomorphic to the direct sum 
$\bigoplus_{\chi\in \widehat G/\Gamma}(\FF^{\ell_\chi},\rho_\chi)$,  
where $\widehat G/\Gamma$ stands for a complete list of representatives of the $\Gamma$-orbits in $\widehat G$. 
\item[(iv)]  $\{(\FF^{\ell_\chi},\rho_\chi)\mid \chi\in \widehat G/\Gamma\}$  is a complete irredundant list of 
representatives of the isomorphism classes of the irreducible representations of $G$ over $\FF$. 
\end{itemize}  
\end{proposition}

\begin{proof} 
(i): Since $A_\chi\in \field^{\ell_\chi\times\ell_\chi}$ has $\ell_\chi$ distinct eigenvalues in $\field$, and the set of eigenvalues 
is the $\Gamma$-orbit of $\omega_\chi$, the space $\field^{\ell_\chi}$ decomposes as the direct sum of $1$-dimensional eigenspaces 
$\field^{\ell_\chi}=\bigoplus_{\omega\in \Gamma\cdot \omega_\chi}\field v_\omega$, 
where $A_\chi v_\omega=\omega v_\omega$, implying that 
$g_\chi \cdot v_\omega =\rho_\chi(g_\chi)v_\omega=A_\chi v_\omega=\omega v_\omega$.  
Moreover, $\ker(\chi)=\ker(\rho_\chi)$, and since the coset of $g_{\chi}$ generates $G/\ker(\rho_\chi)$, 
we conclude that the subspace $\field v_\omega$ is $G$-invariant. 
Consequently, $(\field v_\omega,\rho_\chi\vert_{\field v_\omega})\cong (\field,\psi)$ for some 
$\psi\in \widehat G$. We claim that $\psi=\gamma\cdot \chi$, where $\gamma\in \Gamma$ satisfies 
$\omega=\gamma(\omega_\chi)$. 
Indeed, we have 
\[(\gamma\cdot\chi)(g_\chi)v_\omega=\gamma(\chi(g_\chi))v_\omega
=\gamma(\omega_\chi)v_\omega
=\omega v_\omega=A_\chi v_\omega=g_\chi\cdot v_\omega=\psi(g_\chi)v_\omega,\] 
showing that the characters $\gamma\cdot \chi$ and $\psi$ agree on $g_\chi$. 
Since  $\ker(\psi)\supseteq \ker(\rho_\chi)=\ker(\chi)$,  
$\ker(\gamma\cdot \chi)=\ker(\chi)$, and the coset of $g_\chi$ generates $G/\ker(\chi)$, 
we conclude that $\gamma\cdot \chi=\psi$.  
This shows that $(\field v_\omega,\rho_\chi\vert_{\field v_\omega})\cong (\field,\gamma\cdot\chi)$, 
and thus (i) holds. 

  
(ii):  Let $V$ be a non-zero $G$-invariant subspace in $\FF^{\ell_\chi}$ (where $G$ acts on $\FF^{\ell_\chi}$ via $\rho_\chi$). 
The non-zero $A_\chi$-invariant subspace $\field\otimes_\FF V$  of $\field^{\ell_\chi}$ 
contains a non-zero $A_\chi$-eigenvector $v$. Its  
eigenvalue $\omega$ belongs to $\Gamma\cdot \omega_\chi$. 
The action of $\Gamma$ on $\field$ induces an action on 
$\field^{\ell_\chi}$ and on $\field^{\ell_\chi\times \ell_\chi}$ in the obvious way. The matrix $A_\chi$ is fixed by $\Gamma$, because its entries 
belong to $\FF$. For $\gamma\in\Gamma$ we have 
\begin{equation}\label{eq:gamma v}
A_\chi\gamma(v)=\gamma(A_\chi)\gamma(v)=
\gamma(A_\chi v)=
\gamma(\omega v)=\gamma(\omega)\gamma(v). 
\end{equation}   
Since $V$ is defined over $\FF$, 
$\Gamma$ maps $\field\otimes_\FF V$ into itself, so $\gamma(v)$ belongs to 
$\field\otimes_\FF V$, and it is an $A_\chi$-eigenvector with eigenvalue $\gamma(\omega)$. 
This holds for all $\gamma\in \Gamma$, so $\field\otimes_\FF V$ contains all the $A_\chi$-eigenvectors in 
$\field^{\ell_\chi}$.  
Therefore 
$\field\otimes_\FF V=\field^{\ell_\chi}$, and consequently, $V=\FF^{\ell_\chi}$.   
Thus $\FF^{\ell_\chi}$ has no proper non-zero $G$-invariant subspace, 
so the representation $(\FF^{\ell_\chi},\rho_\chi)$ is irreducible, i.e. (ii) holds. 

(iii)  By (i) we see that 
$\bigoplus_{\chi\in \widehat G/\Gamma}(\field^{\ell_\chi},\rho_\chi)
\cong \bigoplus_{\psi\in \widehat G}(\field,\psi)\cong (\mathcal{F}(G,\field),\rho_{\mathrm{reg}})$. 
So extending the base field from $\FF$ to $\field$ the representations 
$(\mathcal{F}(G,\FF),\rho_{\mathrm{reg}})$ and $\bigoplus_{\chi\in \widehat G/\Gamma}(\FF^{\ell_\chi},\rho_\chi)$ become isomorphic. 
It follows that these two representations are isomorphic over $\FF$. 

(iv) Statement (i) implies that for $\chi\neq \chi'$ in $\widehat G/\Gamma$, the representations 
$(\FF^{\ell_\chi},\rho_\chi)$ and $(\FF^{\ell_{\chi'}},\rho_{\chi'})$ are non-isomorphic, so the given list of irreducible representations is irredundant. 
On the other hand, 
every irreducible representation of $G$ over $\FF$ occurs as a direct summand 
in the regular representation, hence (iii) implies completeness of the list. 
 \end{proof} 

\end{document}
